\BourbakiExerciseGroup

\begin{enumerate}
\def\labelenumi{\arabic{enumi})}
\tightlist
\item
  Let \(r_k\) be the polynomial with integer coefficients such that
\end{enumerate}

\[
p _ {k} = r _ {k} (s _ {1}, \dots , s _ {k}).
\]

\begin{enumerate}
\def\labelenumi{\alph{enumi})}
\item
  Show that \(r_k(\mathbf{Y}_1, \ldots, \mathbf{Y}_k)\) contains the terms \(\mathbf{Y}_1^k\) and \((-1)^{k-1} k\mathbf{Y}_k\) .
\item
  Show that
\end{enumerate}

\[
r _ {k} \left(\mathrm{Y} _ {1}, \dots , \mathrm{Y} _ {k}\right) = \sum_ {p = 0} ^ {k - 1} (- 1) ^ {p} \quad \left| \begin{array}{c c c c c} \mathrm{Y} _ {p + 1} & \mathrm{Y} _ {p + 2} & \mathrm{Y} _ {p + 3} & \dots \mathrm{Y} _ {k - 1} & \mathrm{Y} _ {k} \\ 1 & \mathrm{Y} _ {1} & \mathrm{Y} _ {2} & \dots \mathrm{Y} _ {k - p - 2} & \mathrm{Y} _ {k - p - 1} \\ 0 & 1 & \mathrm{Y} _ {1} & \dots \mathrm{Y} _ {k - p - 3} & \mathrm{Y} _ {k - p - 2} \\ 0 & 0 & 0 & \dots \mathrm{Y} _ {1} & \mathrm{Y} _ {2} \\ 0 & 0 & 0 & \dots 1 & \mathrm{Y} _ {1} \end{array} \right|
\]

\begin{enumerate}
\def\labelenumi{\arabic{enumi})}
\setcounter{enumi}{1}
\item
  Show that \(n! s_n \equiv s_1^n\) modulo the ideal of \(A[X_1, \ldots, X_n]\) generated by \(X_1^2, X_2^2, \ldots, X_n^2\) .
\item
  Let \(f\) be a symmetric polynomial of \(\mathbf{K}[\mathbf{X}_1, \ldots, \mathbf{X}_n]\) and \(\varphi\) the unique polynomial of \(\mathbf{K}[\mathbf{Y}_1, \ldots, \mathbf{Y}_n]\) such that \(f(\mathbf{X}_1, \ldots, \mathbf{X}_n) = \varphi(s_1, \ldots, s_n)\) . Show that the total degree of \(\varphi\) is equal to the degree of \(f\) in any one of the \(\mathbf{X}_i\) .
\end{enumerate}

* 4) Let K be a field of characteristic 0. Show that if \(\alpha_{i}\) ( \(1 \leqslant i \leqslant n\) ) are \(n\) elements of K such that \(\alpha_{1}^{k} + \alpha_{2}^{k} + \cdots + \alpha_{n}^{k} = 0\) for \(n\) consecutive values \(h, h+1, \ldots, h+n-1\) of \(k\) , then \(\alpha_{1} = \alpha_{2} = \cdots = \alpha_{n} = 0\) . Show by an example that the property no longer holds if the \(n\) values of \(k\) are not consecutive or if K is not of characteristic 0.

\begin{enumerate}
\def\labelenumi{\arabic{enumi})}
\setcounter{enumi}{4}
\tightlist
\item
  Let K be a commutative field and \(f \in K[X_{1}, \ldots, X_{n}, Y_{1}, \ldots, Y_{n}]\) ; for every permutation \(\sigma \in S_{n}\) denote by \(\sigma f\) the rational fraction
\end{enumerate}

\[
\sigma f \left(\mathrm{X} _ {1}, \dots , \mathrm{X} _ {n}, \mathrm{Y} _ {1}, \dots , \mathrm{Y} _ {n}\right) = f \left(\mathrm{X} _ {\sigma (1)}, \dots , \mathrm{X} _ {\sigma (n)}, \mathrm{Y} _ {\sigma (1)}, \dots , \mathrm{Y} _ {\sigma (n)}\right).
\]

Then \(f\) is said to be symmetric with respect to the pairs \((\mathbf{X}_i, \mathbf{Y}_i)\) if \(\sigma f = f\) for all \(\sigma \in \mathfrak{S}_n\) .

\begin{enumerate}
\def\labelenumi{\alph{enumi})}
\tightlist
\item
  For every element \(\nu = (\lambda_1, \ldots, \lambda_n, \mu_1, \ldots, \mu_n)\) of \(\mathbf{N}^{2n}\) and every permutation \(\sigma \in \mathfrak{S}_n\) write
\end{enumerate}

\[
\sigma^ {- 1} (\nu) = \left(\lambda_ {\sigma (1)}, \dots , \lambda_ {\sigma (n)}, \mu_ {\sigma (1)}, \dots , \mu_ {\sigma (n)}\right),
\]

so that \(\mathfrak{S}_n\) in this way becomes a group of operators on \(\mathbf{N}^{2n}\) . Let \(\gamma\) be any orbit of \(\mathfrak{S}_n\) in \(\mathbf{N}^{2n}\) , and denote by \(u_{\gamma}\) the symmetric polynomial

\[
\sum \mathrm{X} _ {1} ^ {\lambda_ {1}} \mathrm{X} _ {2} ^ {\lambda_ {2}} \dots \mathrm{X} _ {n} ^ {\lambda_ {n}} \mathrm{Y} _ {1} ^ {\mu_ {1}} \dots \mathrm{Y} _ {n} ^ {\mu_ {n}}
\]

where \((\lambda_{1},\ldots,\lambda_{n},\mu_{1},\ldots,\mu_{n})\) ranges over \(\gamma\) . Show that the polynomials \(u_{\gamma}\) form a basis over K of the vector space of all polynomials that are symmetric in the \((X_{i},Y_{i})\) .

* b) Suppose that K has characteristic 0. Show that every polynomial \(u_{\gamma}\) is equal to a polynomial with rational coefficients in the symmetric functions

\[
\boldsymbol {v} _ {\lambda \mu} = \sum_ {i = 1} ^ {n} \mathbf {X} _ {i} ^ {\lambda} \mathbf {Y} _ {i} ^ {\mu}
\]

(argue by induction on the number of pairs \((\lambda_i, \mu_i)\) which are not \((0, 0)\) in \(\nu = (\lambda_1, \ldots, \lambda_n, \mu_1, \ldots, \mu_n)\) in considering the products \(u_\gamma v_{\lambda \mu}\) ).

\begin{enumerate}
\def\labelenumi{\alph{enumi})}
\setcounter{enumi}{2}
\tightlist
\item
  By elementary symmetric functions of the \((\mathbf{X}_i, \mathbf{Y}_i)\) we shall understand the \(n(n + 3)/2\) polynomials \(w_{hk} = u_\gamma\) corresponding to the orbits \(\gamma\) of elements \((\lambda_1, \ldots, \lambda_n, \mu_1, \ldots, \mu_n)\) other than \((0, \ldots, 0)\) and such that \(h\) of the pairs \((\lambda_i, \mu_i)\) are equal to \((1, 0)\) , \(k\) others are \((0, 1)\) and the remaining \(n - h - k\) are \((0, 0)\) . Show that if \(\mathbf{K}\) is of characteristic 0, then every symmetric polynomial in the \((\mathbf{X}_i, \mathbf{Y}_i)\) is equal to a polynomial in the elementary symmetric functions with coefficients in \(\mathbf{K}\) (consider the sums \(\sum_{i=1}^{n} (\mathrm{UX}_i + \mathrm{VY}_i)^k\) , where \(\mathbf{U}\) and \(\mathbf{V}\) are two indeterminates, and use \(b\) )).
\end{enumerate}

\(d)\) If \(\mathbf{K}\) is a field of characteristic 3, and \(n \geqslant 4\) , show that \(\sum_{i=1}^{n} \mathbf{X}_i^2 \mathbf{Y}_i^2\) is not equal to any polynomial in the symmetric functions \(w_{hk}\) with coefficients in \(\mathbf{K}\) .

\begin{enumerate}
\def\labelenumi{\arabic{enumi})}
\setcounter{enumi}{5}
\tightlist
\item
  Let \(n\) be an integer \(\geqslant 1\) , and let \(X_1, \ldots, X_n\) be indeterminates. For \(1 \leqslant j \leqslant n\) let \(s_j\) be the elementary symmetric polynomial of degree \(j\) in the \(X_i\) and \(u_j = (-1)^{j+1} s_j\) . Put \(S(T) = \prod_{i=1}^{n} (1 - X_i T) = 1 - \sum_{j=1}^{n} u_j T^j\) and for each \(k \geqslant 1\) put \(p_k = \sum_{i=1}^{n} X_i^k\) ; show that
\end{enumerate}

\[
p _ {k} = \sum a _ {\lambda_ {1} \lambda_ {2} \dots \lambda_ {n}} u _ {1} ^ {\lambda_ {1}} u _ {2} ^ {\lambda_ {2}} \dots u _ {n} ^ {\lambda_ {n}}
\]

the summation being extended to all systems \((\lambda_{1},\ldots,\lambda_{n})\) of integers \(\geqslant0\) such that \(\lambda_{1}+2\lambda_{2}+\cdots+n\lambda_{n}=k\) , and the coefficient \(a_{\lambda_{1}\lambda_{2}\ldots\lambda_{n}}\) being equal to

\[
\frac {k \left(\lambda_ {1} + \lambda_ {2} + \cdots + \lambda_ {n} - 1\right) !}{\lambda_ {1} ! \lambda_ {2} ! \dots \lambda_ {n} !}
\]

(expand the series \(S'(T) / S(T)\) in a formal power series in \(T\) ).

\begin{enumerate}
\def\labelenumi{\arabic{enumi})}
\setcounter{enumi}{6}
\tightlist
\item
  With the notation of Exercise 6 show that
\end{enumerate}

\[
\sum a _ {\lambda_ {1} \lambda_ {2} \dots \lambda_ {n}} = - 1 + k \sum_ {j = 0} ^ {k / (n + 1)} \frac {(- 1) ^ {j}}{k - j n} \binom {k - j n} {j} 2 ^ {k - j (n + 1)}.
\]

\begin{enumerate}
\def\labelenumi{\arabic{enumi})}
\setcounter{enumi}{7}
\tightlist
\item
  Put
\end{enumerate}

\[
s _ {k} = \sum \mathrm{A} _ {m _ {1} \dots m _ {n}} p _ {1} ^ {m _ {1}} \dots p _ {n} ^ {m _ {n}},
\]

where the \(\mathbf{A}_{m_1\dots m_n}\) are rational numbers and the sum is extended over all families \((m_1,\dots,m_n)\) such that \(\sum i\cdot m_i = k\) . Show that \(\sum \mathbf{A}_{m_1\dots m_n} = 0\) for \(k > 1\) and \(\sum |\mathbf{A}_{m_1\dots m_n}| = 1\) for \(k\geqslant 1\) . Deduce that on writing \(q_{j} = -p_{j}\) we have

\[
(- 1) ^ {k} s _ {k} = \sum \left| \mathrm{A} _ {m _ {1} \dots m _ {n}} \right| q _ {1} ^ {m _ {1}} \dots q _ {n} ^ {m _ {n}}.
\]

\begin{enumerate}
\def\labelenumi{\arabic{enumi})}
\setcounter{enumi}{8}
\tightlist
\item
  \begin{enumerate}
  \def\labelenumii{\alph{enumii})}
  \tightlist
  \item
    Let \(\mathbf{K}\) be a field and \(\mathbf{X}_1, \ldots, \mathbf{X}_n\) indeterminates. Show that for any exponents \(\nu_1 < \nu_2 < \ldots < \nu_n\) the determinant \(\det(\mathbf{X}_i^{\nu_j})\) is non-zero in \(\mathbf{K}(\mathbf{X}_1, \ldots, \mathbf{X}_n)\) .
  \end{enumerate}
\end{enumerate}

\begin{enumerate}
\def\labelenumi{\alph{enumi})}
\setcounter{enumi}{1}
\tightlist
\item
  Let \(s_k\) ( \(1 \leqslant k \leqslant n\) ) be the elementary symmetric polynomials in the \(X_j\) and \(r_1, \ldots, r_n\) \(n\) elements of the field \(K(X_1, \ldots, X_n)\) . We put (with a new indeterminate T)
\end{enumerate}

\[
\mathrm{R} (\mathrm{T}) = \frac {1 + r _ {1} \mathrm{T} + \cdots + r _ {n} \mathrm{T} ^ {n}}{1 + s _ {1} \mathrm{T} + \cdots + s _ {n} \mathrm{T} ^ {n}} = 1 + u _ {1} \mathrm{T} + \dots + u _ {k} \mathrm{T} ^ {k} + \dots .
\]

Show that if \(n\) of the coefficients \(u_{k}\) are zero, then all the \(u_{k}\) are zero. (Decompose \(R(T)\) into simple fractions.)

* \(c\) ) Suppose that \(\mathbf{K}\) is of characteristic zero. Let \(\{\nu_1,\nu_2,\dots,\nu_n\}\) be a set of \(n\) elements \(\geqslant 1\) of \(\mathbf{N}\) , and let \(\mathbf{M}\) be the complement of this set in \(\mathbf{N}\) ; suppose also that \(\mathbf{M}\) is stable under addition. Show that if \(y_{1},\ldots ,y_{n}\) are \(n\) elements of \(\mathbf{K}(\mathbf{X}_1,\dots,\mathbf{X}_n)\) such that

\[
p _ {\nu _ {j}} (y _ {1}, \dots , y _ {n}) = p _ {\nu _ {j}} (\mathbf {X} _ {1}, \dots , \mathbf {X} _ {n})
\]

for \(1 \leqslant j \leqslant n\) , then \(y_{j} = X_{j}\) for \(1 \leqslant j \leqslant n\) except possibly for a permutation. (Consider the expansion of \(\log R(T) = v_{1}T + v_{2}T^{2} + \cdots + v_{k}T^{k} + \cdots\) , observe that \(v_{\nu_{j}} = 0\) for \(1 \leqslant j \leqslant n\) and examine the relations between the \(v_{k}\) and the \(u_{k}\) .) *

\begin{enumerate}
\def\labelenumi{\arabic{enumi})}
\setcounter{enumi}{9}
\item
  Let \(a \in \mathbf{A}\) . For each polynomial \(h \in \mathbf{A}[\mathbf{X}]\) put \(h_a(\mathbf{X}) = h(\mathbf{X} + a)\) ; show that \(\operatorname{res}(f_a, g_a) = \operatorname{res}(f, g)\) , \(\operatorname{dis}(f_a) = \operatorname{dis}(f)\) .
\item
  In this exercise we assume that A is a field, \(f(0) \neq 0\) and \(\deg g < \deg f = p\) . Let \(\sum \alpha_{n} X^{n}\) be the expansion in a formal power series of the rational fraction \(g(X) / f(X)\) ; for each integer \(k \geqslant 0\) let \(H_{0}^{(k)}\) be the Hankel determinant defined in Exercise 1 of IV, p.~90. a) Show that \(H_{0}^{(k)} = 0\) for \(k > p\) and calculate \(H_{0}^{(p)}\) as function of \(\operatorname{res}(f, g)\) .
\end{enumerate}

\begin{enumerate}
\def\labelenumi{\alph{enumi})}
\setcounter{enumi}{1}
\tightlist
\item
  Under what conditions do we have \(\mathbf{H}_0^{(p)} = \mathbf{H}_0^{(p - 1)} = \ldots = \mathbf{H}_0^{(p - r + 1)} = 0\) ?
\end{enumerate}

\begin{enumerate}
\def\labelenumi{\arabic{enumi})}
\setcounter{enumi}{11}
\tightlist
\item
  There exists a unique polynomial \(P_{p,q}(T_{0},\ldots,T_{p},U_{0},\ldots,U_{q})\) with integer coefficients such that for every ring A and any polynomials \(f=t_{p}X^{p}+\cdots+t_{0}, g=u_{q}X^{q}+\cdots+u_{0}\) of A{[}X{]} we have
\end{enumerate}

\[
\operatorname{res} _ {p, q} (f, g) = \mathrm{P} _ {p, q} \left(t _ {0}, \dots , t _ {p}, u _ {0}, \dots , u _ {q}\right).
\]

If we assign the weight i to \(T_{i}\) and \(U_{i}\) , then \(P_{p,q}\) is isobaric of weight pq.

\begin{enumerate}
\def\labelenumi{\arabic{enumi})}
\setcounter{enumi}{12}
\tightlist
\item
  Let \(n \geqslant 1\) be an integer, and order the elements of \(\mathbf{N}^n\) lexicographically (notation \(\alpha \preccurlyeq \beta\) ). For each \(\alpha \in \mathbf{N}^n\) let \(T_{\alpha}\) be the set of elements of \(\mathbf{Z}[X_1, \ldots, X_n]\) of the form \(X^{\alpha} + \sum_{\beta < \alpha} u_{\alpha\beta} X^{\beta}\) (with \(u_{\alpha\beta} \in Z\) ).
\end{enumerate}

\begin{enumerate}
\def\labelenumi{\alph{enumi})}
\tightlist
\item
  Show that the product of an element of \(\mathbf{T}_{\alpha}\) by an element of \(\mathbf{T}_{\beta}\) lies in \(\mathbf{T}_{\alpha + \beta}\) .\\
\item
  Let \(\varepsilon_{k}\) be the element \((1, 1, \ldots, 1, \underbrace{0, \ldots, 0}_{k})\) of \(\mathbf{N}^{n}\) . Show that \(s_{k} \in \mathbf{T}_{\varepsilon_{k}}\) . Deduce that
\end{enumerate}

\(S(\alpha) = \prod_{k=1}^{n} s_k^{\alpha_k - \alpha_{k+1}} \text{ belongs to } T_\alpha, \text{ for } \alpha \in \mathbb{N}^n \text{ such that } \alpha_1 \geqslant \alpha_2 \geqslant \ldots \geqslant \alpha_n. (\text{ Set } \alpha_{n+1} = 0.)\) \(c) \text{ With the notation of IV, p. 66, deduce that } c_{\alpha\alpha} = d_{\alpha\alpha} = 1 \text{ for } \alpha \in D_k \text{ and } c_{\alpha\beta} = d_{\alpha\beta} = 0 \text{ if } \alpha < \beta.\)

\begin{enumerate}
\def\labelenumi{\arabic{enumi})}
\setcounter{enumi}{13}
\tightlist
\item
  Let \(L\) be a finite ordered set and \(A\) a commutative ring. In the ring \(M\) of square matrices over \(A\) with index set \(L\) , the set of matrices \(m = (m_{uv})\) such that \(m_{uv} = 0\) except when \(u \leqslant v\) is a subring \(B\) .
\end{enumerate}

\begin{enumerate}
\def\labelenumi{\alph{enumi})}
\item
  An element \(m = (m_{uv})\) of B is invertible in B if and only if \(m_{uu}\) is invertible in A for all \(u \in L\) .
\item
  Suppose henceforth that \(\mathbf{A} = \mathbf{Z}\) . Let \(\zeta\) be the element of \(\mathbf{B}\) defined by \(\zeta(u, v) = 1\) if \(u \leqslant v\) and \(\zeta(u, v) = 0\) otherwise. Show that \(\zeta\) has an inverse \(\mu\) (also written \(\mu_{\mathrm{L}}\) ) in \(\mathbf{B}\) (called the Möbius function of L). Let \(f\) and \(g\) be two mappings of L into \(\mathbf{Z}\) . Show that the relations \(\ll f(x) = \sum_{y \geqslant x} g(y)\) for all \(x \in L\) and \(\ll g(x) = \sum_{y \geqslant x} \mu(x, y) f(y)\) for all \(x \in L\) are equivalent.
\item
  Assume that \(L\) is the product of a finite family of ordered sets \(L_1, \ldots, L_n\) . Show that \(\mu_L(x, y) = \prod_{i=1}^{n} \mu_{L_i}(x_i, y_i)\) for \(x = (x_1, \ldots, x_n)\) and \(y = (y_1, \ldots, y_n)\) such that \(x \leqslant y\) .
\item
  Let \(m \geqslant 1\) be an integer and let \(D\) be the set of divisors of \(m\) , equipped with the order relation « \(x\) divides \(y\) ». Show that \(\mu_{D}(x, y) = (-1)^{s}\) where \(s\) is the cardinal of the set of prime numbers \(p\) dividing \(y/x\) when \(y/x\) is not divisible by the square of an integer, and \(\mu_{D}(x, y) = 0\) otherwise. (By applying \(c\) ) reduce to the case where \(m\) is a power of a prime number.)
\item
  Suppose that the ordered set L is lattice-ordered. Let \(a, b, c\) be distinct elements of L such that \(a \leqslant b \leqslant c\) , and let \(\Phi\) be the set of \(x \in \mathbf{L}\) such that \(a \leqslant x\) and \(\sup(x, b) = c\) , and \(\Psi\) the corresponding set defined by \(y \leqslant c\) and \(\inf(y, b) = a\) . Show that we have the relation \(\sum_{x \in \Phi} \mu(a, x) = \sum_{y \in \Psi} \mu(y, c) = 0\) .
\item
  Let T be a finite set and \(\mathcal{R}(T)\) the set of equivalence relations on T, ordered by the relation « R is finer than S », written \(R \leqslant S\) . Show that \(\mathcal{R}(T)\) is lattice-ordered. Let R, S in \(\mathcal{R}\) be such that \(R \leqslant S\) ; let \(U_1, \ldots, U_k\) be the equivalence classes under S and for \(1 \leqslant i \leqslant k\) , let \(\alpha_i\) be the number of equivalence classes under R contained in \(U_i\) . Show that \(\mu_{\mathcal{R}(T)}(R, S) = \prod_{i=1}^{k} (-1)^{\alpha_i - 1} (\alpha_i - 1)!\) (if \(A_i\) is the set of equivalence classes under R
\end{enumerate}

contained in \(U_{i}\) , observe that the interval \([R, S]\) of \(\mathcal{R}(T)\) is isomorphic to \(\mathcal{R}(A_{1}) \times \ldots \times \mathcal{R}(A_{k})\) , apply c) and reduce to the case where R is the least element of \(\mathcal{R}(T)\) and S the greatest; treat this last case by using e)).

\begin{enumerate}
\def\labelenumi{\arabic{enumi})}
\setcounter{enumi}{14}
\tightlist
\item
  Let \(\mathbf{T}\) be a finite set with \(k\) elements. We denote by \(\mathbf{D}_k\) the set of elements \((\alpha_1, \ldots, \alpha_k)\) of \(\mathbb{N}^k\) such that \(\alpha_1 + \ldots + \alpha_k = k\) and \(\alpha_1 \geqslant \alpha_2 \geqslant \ldots \geqslant \alpha_k\) . With every equivalence relation R in T we associate an element \(\lambda(R)\) of \(D_k\) as follows: if \(U_1, \ldots, U_l\) are the equivalence classes under R, numbered so that \(\operatorname{Card}(U_1) \geqslant \ldots \geqslant \operatorname{Card}(U_l)\) , put \(\lambda(R) = (\alpha_1, \ldots, \alpha_k)\) with \(\alpha_i = \operatorname{Card}(U_i)\) if \(1 \leqslant i \leqslant l\) and \(\alpha_i = 0\) if \(l < i \leqslant k\) . Let \(X_n (n \in N)\) be indeterminates; for each \(R \in \mathcal{R}(T)\) put \(r(R) = \sum_{f} \prod_{t \in T} X_{f(t)}\) , where the summation is extended to the set of mappings \(f: T \to N\) such that R is the equivalence relation associated with \(f\) . Put \(s(R) = \sum_{S \geqslant R} r(S)\) .
\end{enumerate}

\begin{enumerate}
\def\labelenumi{\alph{enumi})}
\tightlist
\item
  Prove that \(r(\mathbf{R}) = \prod_{i=1}^{k} r_i \cdot \mathbf{M}(\lambda(\mathbf{R}))\) if \(r_i\) of the equivalence classes under \(\mathbf{R}\) are of
\end{enumerate}

cardinal \(i\) for \(1 \leqslant i \leqslant k\) . (See formula (11) of IV, p.~65 for the definition of \(\mathbf{M}(\alpha)\) .)

\begin{enumerate}
\def\labelenumi{\alph{enumi})}
\setcounter{enumi}{1}
\item
  Prove the relation \(S(R) = \prod_{i=1}^{k} p_i^{r_i}\) where the \(r_i\) are as in a) and \(p_m = \sum_{n \in N} X_n^m\) for every integer \(m \geqslant 1\) .
\item
  Show that \(r(\mathbf{R}) = \sum_{\mathbf{S} \geqslant \mathbf{R}} \mu(\mathbf{R}, \mathbf{S}) s(\mathbf{S})\) , where \(\mu(\mathbf{R}, \mathbf{S})\) is defined as in Exercise 14.
\item
  Hence deduce an expression of the \(s_k\) as a polynomial in \(p_1, \ldots, p_k\) with rational coefficients, where \(s_k = \sum_{i_1 \leq \ldots \leq i_k} X_{i_1} \ldots X_{i_k}\) .
\item
  For \(\alpha \in \mathbf{D}_k\) put \(S'(\alpha) = \prod_{i=1}^{k} s_{\alpha_i}\) . Show that \(S'(\alpha) = \sum_{\beta \in D_k} c_{\alpha\beta}' \cdot M(\beta)\) , where the integer \(c_{\alpha\beta}'\) is defined as follows: suppose that \(r_i\) (resp. \(s_i\) ) of the components of \(\alpha\) (resp. \(\beta\) ) are equal to \(i\) (for \(1 \leqslant i \leqslant k\) ). Then \(c_{\alpha\beta}'\) is the product \(Nr_1!... r_k!s_1!... s_k!/k!\) , where N is the number of pairs of equivalence relations R, S on T such that \(\lambda(R) = \alpha, \lambda(S) = \beta\) and such that the intersection of an equivalence class under R with an equivalence class under S has at most one element. Deduce the symmetry relation \(c_{\alpha\beta}' = c_{\beta\alpha}'\) .
\item
  Show that there is an involution \(\alpha \mapsto \tilde{\alpha}\) in \(D_k\) , characterized by the fact that the relations \(i \leqslant \alpha_j\) and \(j \leqslant \tilde{\alpha}_i\) are equivalent, for \(i\) , \(j\) in the range from 1 to \(k\) . Show that \(S'(\alpha) = S(\tilde{\alpha})\) and \(c_{\alpha\beta}' = c_{\tilde{\alpha}\beta}\) . Hence obtain the relation \(c_{\alpha\beta} = c_{\tilde{\alpha}\tilde{\beta}}\) (with the notation of IV, p.~66).
\end{enumerate}

